\subsection{Frozen Empirical Spectrum Predictions}\label{sec:spectralpredictionmethods}

We test how much of the two-round spectrum can be predicted from positional HW distributions and independently sampled residuals. The models are fitted only to the earlier 80-key-pair, 8192-bit validation cohort (indices 840--919), using its two-round observations. The evaluation uses 80 new key pairs (1200--1279), all 64 singleton nonce masks, four shared base nonces and 8192 output bits. Standard additive output is used for both ciphers. The model specifications, Monte Carlo budget and reporting rules were fixed before this evaluation; fitted parameters and predictions were hashed before generating any evaluation streams.

Model M0 pools an empirical HW32 distribution for each cipher, nonce mask and output word across training keys, bases and blocks, then samples every new word independently. Model M1 instead draws one complete training-key index per simulated key pair, shared across both ciphers and all masks, and retains its four base-nonce slots. Within each key/mask/base/word layer it independently samples the empirical distribution over the 16 training blocks. M1 therefore retains conditional positional profiles and their shared variation across masks while removing residual dependence between words and blocks. Its layers contain only 16 observed blocks; repeated values are not independent training replicates. Both models preserve observed support zeros, including those implied by the two-round identities. The remaining fitted distributions are empirical rather than consequences of those identities.

For segment width $s$, constituent word means and model variances are summed, and the mean profile is repeated over 16 blocks. If $\ell$ denotes a model layer, independent residuals give
\begin{equation}
 \mathbb{E}\!\left[\left|\operatorname{DFT}(h^{(\ell)})_k\right|^2\mid\ell\right]
 =\left|\operatorname{DFT}(\mu^{(\ell)})_k\right|^2
  +\sum_{t=0}^{T-1}\sigma_{t,\ell}^{2},\qquad T=8192/s.
 \label{eq:predictedpower}
\end{equation}
For M1, the four base-conditional powers are averaged within a selected key, then over training keys; applying a Fourier transform to the pooled mean would discard between-layer variation. The entropy of normalized expected power is not expected finite-sample entropy. We therefore generate 4096 complete pseudo-key pairs per model, pool $Q=4$ Fourier powers within each mask, and apply the same finite-sample diagnostics as to the evaluation data. One HW32 draw is aggregated to all three widths, 32, 64 and 128 bits. Independent $\operatorname{Binomial}(32,1/2)$ word weights provide the finite-model reference; the actual evaluation uses input-reuse-matched SHAKE256 streams.

All six cipher-by-width cells are reported for both models: predicted and observed $G$, reference-adjusted 512-bit fundamental share, and normalized spectra. Signed errors are evaluation means minus frozen predictions. Pointwise intervals use 5000 whole-evaluation-key resamples, conditional on the fixed training fit and Monte Carlo estimate; Monte Carlo standard errors are reported separately. We also report the $L^1$ distance between predicted and observed mean AC-normalized spectra. This averaged distance can conceal compensating errors of individual masks or keys, and a positive lower bootstrap bound is not a model-rejection test. No success threshold, hypothesis test, detector calibration or post-evaluation model selection is applied. The comparison tests these two empirical model constructions, not the full independence structure of the cipher.
